\input{determinant/preamble}
\title{Toward a conventional proof of Dirichlet GRH\\
\large Short forms, spanning, and a determinant}
\author{GRH distillation project}
\date{November 2026}
\begin{document}
\maketitle
\begin{abstract}
We reconstruct the Dirichlet-GRH argument through one arithmetic
determinant comparison and its analytic prime-sum endpoint.
The analytic implication and the abstract spanning, determinant
and simultaneous small-twist theorems are proved under their
stated hypotheses. The common arithmetic realization and its
uniform estimates remain open. This is the active conventional
manuscript, not a proof claim based on external formal verification.
\end{abstract}
\tableofcontents
\input{determinant/body}
\section{All moduli and the principal pole}
Given any positive modulus \(N_0\), choose a fixed multiple
\(q\) of \(\operatorname{lcm}(18480,N_0)\).
Inducing the character to modulus \(q\) adds only factors
\(1-\chi(p)p^{-s}\) nonzero in the open critical strip.
Consequently (1) for all admissible \(q,d\) would give the
critical-line assertion for all Dirichlet characters. At \(s=1\)
the principal factor is treated by pole removal; the functional
equation is applied to the underlying primitive character.
The conclusion is conditional on the arithmetic suppliers
in Sections 3--5. It is not general Hecke GRH.
\end{document}
